\section{Introduction}\label{sec:intro}

Virtualization based on containers has gained a lot of interest during last few years. Containers have become popular since they offer a lightweight system abstraction which includes sufficient components for running stand-alone and isolated cloud applications.

Despite the attention paid to security issues in the cloud computing environment, containers have been demonstrated vulnerable to several attacks \cite{mattetti2015securing,abed2015intrusion}. In general scenarios, a container is considered healthy at the start of its operation, and later it is infected with a malware. But in some cases, the container images which are stored in repositories and made available to the community (e.g., dockerhub), might be already infected with a malware. The malware can then be activated silently when the container image is installed and executed. Cryptohijaking malware is one example of malicious software that could be pre-injected in a container image. When executed, the cryptominer hiddenly shares the computation resources with the legitimate application. Additional examples of malicious code that can run inside containers are: Backdoors, Ransomware, External file transfer malware, etc.

According to some recent reports \cite{Backdoored}, several images available on dockerhub which have been downloaded millions of times, have been found with a malicious code pre-installed. Other reports showed that many servers in telecommunications, media, healthcare, hosting containers have been found infected with cryptocurrency miners. Cryptohijaking attacks are the widest threat targeting cloud containers nowadays.

During application deployment and execution in the containers, a malware process can hiddenly start and run in the background, this may lead to resource exhaustion, privacy breaching, and service disruption. Cloud service providers need to detect malicious activities within containers to protect its resources (storage, bandwidth, etc), protect user resources (data, web services, etc), and avoid the spread of the attack from one container to the others. At the same time, service providers cannot implement intrusive anomaly detection strategies which look at the information inside the container's resources.

Intrusion detection is commonly known as a hard and long-standing problem. It becomes more challenging in public cloud computing services when the monitoring system has not direct access to the resources assigned to  customers and their activities within the container processes. In addition, accurate monitoring cloud services is challenging because of the complex architectures and the large-scale of resources to monitor. For this reason, advanced detection solutions can play a crucial role in detecting malicious behaviors by analyzing the data related to the containers and observable at the hosting operating system which share the kernel with the hosted containers. Some features that a machine learning-based detector may exploit for this purpose can be derived by isolating and processing container-based system calls. System calls are the communication instrument between the running application and the Linux kernel of the hosting machine and can be captured by the cloud provider without accessing the container resources directly.
In this paper we provide a mechanism for (i) classifying container workload behaviors and (ii) detecting malicious behaviors within the containers. Our solution, working at the hosting operating system, is based on a mix of supervised and unsupervised machine learning models. While the (supervised) random forest algorithm is able the delineate the differences among the multiple container workloads which may run on the hosting machine, the (unsupervised) isolation forest model will identify anomalous behaviors in the containers. We exploit the anonymous walks embedding algorithm to extract a compact representation of the sequences of system calls which is given as input into the intrusion detection model. The use of graph representation by means of anonymous walks embedding allows to incorporate the hidden relations and dependencies between system calls within a time window.
To the best of our knowledge, we are the first to use anonymous walks for embedding a sequence of system calls into a machine learning model in the context of intrusion detection.

The contributions of this research work are as follow:
\begin{itemize}
  \item We provide a description of a data repository for intrusion detection in the context of cloud containers, which has been collected in a emulated environment running selected Docker workloads;
  \item We introduce a graph representation of the sequence of system calls collected at the hosting machine's kernel level which is processed by the random walks and anonymous walks algorithms to extract features of the container's behaviors;
  \item We present the design of a practical intrusion detection mechanism which incorporates the capability of distinguish among multiple container workloads. Our solution works on three stages and is based on an anonymous walks data embedding stage followed by a random forest classification module and an ensemble of independent isolation forests;
  \item Finally, we provide some results obtained by testing our solution on two datasets: one resulting from a public repository and another obtained from traces collected for three types of workloads.
\end{itemize}

The rest of the paper is organized as follows; We provide the state of the art related to this domain in Section \ref{sec:relworks}. Section \ref{sec:background} provides background information on cloud containers, graph theory, random and anonymous walks algorithms, random forest and isolation forest models, which are required to understand the solution proposed in this work. We describe our graph-based intrusion detection system in Section \ref{sec:model}. The setup for data collection in a cloud environment is described in Section \ref{sec:setup}. Section \ref{sec:results} presents the obtained results, and finally our conclusion is provided in Section \ref{sec:concl}.

